\paragraph{Final model.}
We averaged the positive-class probabilities from the three
RoBERTa-base runs trained with 10\% warmup and linear decay. At a
threshold of 0.5, the population standard deviation of single-run
Macro-F1 decreased from 20.27 percentage points under the original
schedule to 0.43 percentage points across the same three seeds. On the
validation set, the final RoBERTa ensemble achieved 92.73\% Macro-F1
after threshold selection. The 75\% RoBERTa and 25\% TF-IDF + SVM
blend achieved 92.57\% after its own threshold selection, so we used
the RoBERTa ensemble for submission.

\paragraph{Decision threshold.}
On the 1,492 labelled validation reviews, a threshold of 0.590 gave
the highest accuracy (96.25\%) and Macro-F1 (92.73\%) among the
thresholds evaluated for the RoBERTa ensemble. We applied the same
threshold to the test set, assigning label 1 to reviews with an
average positive-class probability of at least 0.590 and label 0
to the remaining reviews.

\paragraph{Test predictions.}
The table shows the predicted label distribution for the 1,851
unlabeled test reviews.

\begin{center}
\begin{tabular}{lrr}
\hline
Predicted label & Number of reviews & Share \\
\hline
Positive (1)    & 1,568             & 84.71\% \\
Negative (0)    &   283             & 15.29\% \\
\hline
Total           & 1,851             & 100.00\% \\
\hline
\end{tabular}
\end{center}

The file \texttt{submission.csv} contains the \texttt{reviews}
and \texttt{predictions} columns in the original test-set order.